\section{Conclusion}
\label{sec:conclusion}

We formulated sparse-PMU event attribution through competing physical
histories. Load/line predictors start from estimated past states and retain
the cause of the response after restoration. Active mass $M$ and its
increment $\Delta M$ expose the changing explanation rather than replacing
the physical dynamics.

\IfFileExists{generated/c01_macros.tex}{The independent continuous
evaluation gives \COneMixTimely{}/\COneMixTotal{} timely family/source
decisions with \COneMixExtras{} extra alarms for the mixture, and
\COneNomTimely{}/\COneNomTotal{} with \COneNomExtras{} for nominal.
Stratified results and paired intervals determine where each choice is
useful.}{In completed development streams, the mixture obtains
\ContMixTimely{}/\ContMixTotal{} timely family/source decisions with
\ContMixExtras{} extra alarms. Nominal gains one timely decision at a
higher alarm burden. Independent confirmation remains in progress.}

\IfFileExists{generated/ei_x02_macros.tex}{The controlled load test distinguishes reference correction from signature conditioning. The mixed term reduces response error and median relative severity error, but does not establish a directional improvement in joint source decisions. This component result does not change the frozen causal method.
}{}

The result is a specified physical attribution procedure and a testable
separation of recovery, reference error and source discrimination. Its
scope remains the declared simulation contracts. It does not establish
all-family diagnosis, universal source recovery, calibrated probabilities
or real-time deployment.
